\chapter{Conclusion}

This thesis set out to answer two questions about LLM-as-a-Judge pipelines for Retrieval-Augmented Generation. First, how badly are contemporary LLM judges fooled by corpus poisoning, and how does that vulnerability vary with attack type, attack severity, and the choice of judge? Second, can a lightweight upstream pre-filter mitigate the failure?

The experimental programme answered both questions on data. The vulnerability is real and substantial: baseline faithfulness False Positive Rates across three contemporary judges range from 0.412 to 0.699, and the three judges occupy distinct calibration regimes that respond differently to attack class. The most adversarial attack class (PoisonedRAG-style) produces a directional cross-judge reversal, where the judges that find this attack hardest are the ones with continuous or extreme-leniency scoring, and the judge that finds it easiest is the one with near-binary scoring. The vulnerability is not uniform across the LLM-as-a-Judge paradigm; it is structured by judge architecture and attack type.

The second answer is more uncomfortable. Pre-filtering, in the two designs we evaluated, does not reliably mitigate the vulnerability. Decomposing originally-poisoned triplets into True (supporting passages overwritten) and Survived (supporting passages intact) reveals that filtering directionally raises True-category faithfulness for all three judges under the statistical filter (Δ +0.053 GPT, +0.072 Gemma, +0.151 DeepSeek). The directional pattern is consistent across judges, but the statistical robustness varies: under question-clustered bootstrap, only DeepSeek's True-subset effect approaches significance (corrected p = 0.099), while GPT's and Gemma's True-subset effects do not survive clustered correction (corrected p = 0.239 and 0.259 respectively). The largest and most robust single end-to-end effect is in the Survived subset for PoisonedRAG-style triplets, where DeepSeek's mean faithfulness rises by 0.377 under the statistical filter (clustered bootstrap p = 0.012). On context relevance the paradox is concentrated on DeepSeek (Δ +0.180); on answer relevance changes are small across all cells.

The operative mechanism is context clarification through distractor removal. The statistical filter, designed to remove poisoned passages, also removes a substantial fraction of the original HotpotQA distractor passages. The residual is shorter (mean 7.1 versus 11.6 passages, a 39\% reduction) and more topically focused, and judges score this cleaner residual as more faithful regardless of whether it actually contains the supporting evidence. We tested the alternative collateral-damage hypothesis (filter removes supporting passages, judge loses contradicting evidence) directly: it is empirically refuted. In triplets where the filter removed supporting passages, scores fell across all three judges by 0.15 to 0.23 - the judges correctly registered the missing evidence. The paradox emerges only in triplets where the filter spared the supporting passages, and is most pronounced where the filter removed the most original distractors.

These two findings - that judges are vulnerable, and that filtering paradoxically worsens the worst case - together support a clear practical conclusion: pre-filtering, in its current form, is not a sufficient defence for LLM-as-a-Judge deployments under adversarial conditions. Practitioners deploying RAG evaluation pipelines in security-critical contexts should not assume that adding upstream filtering improves downstream reliability. The empirical evidence suggests it can do the opposite.

\section{Restating the Contributions}

The four contributions of the thesis can be restated concisely.

The Poisoned Evaluation Dataset (v2) provides 2,400 question-context-answer triplets from HotpotQA across three injection types and four noise levels, with explicit per-passage poisoning indices. The PoisonedRAG-style injection is implemented carefully, specifying a target wrong answer per question and using the full retrieved context. The dataset is released as a reusable benchmark for further work on judge robustness and RAG security.

The multi-metric meta-evaluation framework characterises three contemporary LLM judges across all three RAGAS dimensions under controlled poisoning. The framework reports judge behaviour by injection type, by noise intensity, and by score-distribution shape. It is reusable for future judge-robustness studies and provides the structural basis for the empirical findings reported here.

The multi-signal statistical pre-filter combines five engineered signals through XGBoost into a passage-level classifier. The filter is trained on a stratified split of the Poisoned Evaluation Dataset and evaluated on a held-out test split. It achieves F1 of 0.929 with near-zero clean false-positive rate. A complementary triplet-level Mistral filter is implemented for cross-architecture comparison.

The end-to-end evaluation methodology measures how each pre-filter affects judge reliability. It uses the True/Survived decomposition to distinguish between cases where the attack hit the supporting context and cases where it did not, providing an operationally meaningful measure of post-filter judge failure that is robust to the labelling artefacts that naive analyses suffer from. The methodology surfaces a counterintuitive empirical pattern - filter-induced score increases driven by distractor removal rather than by collateral support loss - that current filter design does not anticipate.

\section{Answers to the Research Questions}

\textbf{RQ1: How badly are contemporary LLM judges fooled by corpus poisoning?}

Substantially. Baseline faithfulness FPRs for the three judges studied range from 0.412 (DeepSeek) to 0.555 (GPT) to 0.699 (Gemma). Context relevance FPRs range from 0.326 to 0.860 and answer relevance FPRs from 0.652 to 0.889. The three judges occupy three qualitatively distinct calibration regimes (near-binary, continuous, extreme-leniency) and respond differently to attack class. PoisonedRAG-style is the hardest attack for GPT and Gemma but the easiest for DeepSeek. Vulnerability decreases monotonically with noise level for all three judges within each injection type at baseline.

\textbf{RQ2: How does that vulnerability vary with judge calibration?}

It varies substantially. The three calibration regimes correspond to three different scoring strategies, each with distinct trade-offs in detection profile, abstention support, and response to filter intervention. Near-binary judges effectively function as contradiction detectors and handle PoisonedRAG-style attacks well; continuous-scoring judges weigh evidence in graded fashion and are more vulnerable to subtle attacks. Extreme-leniency judges have a default-faithful prior that makes them most vulnerable across the board.

\textbf{RQ3: Can a lightweight upstream pre-filter mitigate the failure?}

In the two designs evaluated, no - at least not reliably. The statistical filter is a strong classifier in isolation (F1 = 0.929 on the held-out test split) and removes most poisoned passages. But end-to-end, decomposed by True/Survived, the filter directionally raises the rate at which judges score concentrated-poison contexts as faithful. The directional pattern is consistent across all three judges (+0.053 GPT, +0.072 Gemma, +0.151 DeepSeek mean faithfulness on the True subset), but statistical robustness varies under question-clustered bootstrap: only DeepSeek's True-subset effect approaches significance after correction (corrected p = 0.099); GPT's and Gemma's do not (corrected p = 0.239 and 0.259). The largest single cell-level effect, and the most statistically robust, is on the Survived subset for PoisonedRAG-style attacks, where DeepSeek's mean faithfulness rises by 0.377 (clustered bootstrap p = 0.012). The LLM filter (Mistral) shifts judge response materially only for DeepSeek. On context relevance the paradox is concentrated on DeepSeek; on answer relevance changes are minimal across the board.

\textbf{RQ4: What mechanisms explain the failure of pre-filtering to mitigate?}

The mechanism is now empirically grounded. The filter removes original HotpotQA distractor passages at high rates, producing a shorter and more topically focused residual that judges score as more faithful regardless of content. The collateral-damage hypothesis (filter removes supporting passages, judge loses evidence, paradox emerges) is empirically refuted by a direct counterfactual: when the filter removed supporting passages, scores fell across all three judges, exactly as the alternative hypothesis would predict for an unbiased reasoner. The paradox is driven by the filter's incidental clarification of the residual context, not by the filter's removal of supporting evidence.

\section{Future Work}

Four directions emerge from the present work as the most natural extensions.

\textbf{Cross-task generalisation.} All experiments in this thesis use HotpotQA, a multi-hop QA benchmark with structural properties (multiple cooperating supporting passages, distractor passages, explicit answer supervision) that informed both the attack design and the filter design. Whether the headline paradox - filtering raises True-category faithfulness - generalises to single-hop QA, dialogue evaluation, summarisation faithfulness scoring, or open-domain agentic evaluation is unknown. Replicating the dataset construction, filter design, and True/Survived analysis on at least one structurally different evaluation task is the most direct test of the thesis's central claim. We expect the distractor-removal mechanism to transfer to any task where retrieval delivers a mix of supporting passages and topically loose distractors and where filtering acts on length-shortening principles. Tasks without distractors in the retrieval set (single-hop QA with high-precision retrieval, structured-data evaluation) may not show the same effect at all, since the filter would have less original noise to remove.

\textbf{Judge-side hardening.} The qualitative justification analysis in Section 4.6 (n=28 triplets, single judge) suggests that explicit poison-aware prompting can shift judge behaviour on the engineered attack types (mean faithfulness on PoisonedRAG-style drops from 0.800 to 0.492 on the subset; on adversarial fact from 0.758 to 0.558). The sample is too small to support strong quantitative claims, but the directional shift is large enough to be worth investigating rigorously. We see this as the most empirically promising direction \textit{suggested} by our results, with the caveat that the suggestion rests on a small qualitative subset rather than a full evaluation. Quantifying the effect at scale and characterising its limits across attack types and judges is a tractable research programme.

\textbf{Hybrid filter architectures matched to the True/Survived decomposition.} The statistical filter and the LLM filter detect partially complementary phenomena. A hybrid that combines both would, in principle, achieve higher recall - but the paradox analysis suggests that higher recall may not translate into improved end-to-end judge reliability if the distractor-removal mechanism persists. The relevant filter-design question is whether one can remove poison without also removing original distractors, since the latter is what produces the cleaner-looking residual that judges over-trust. Filter architectures that operate on a different abstraction level (re-ranking rather than removal, contradiction detection that flags inconsistency without modifying content, or filters trained explicitly on the distinction between adversarial poison and naturally-occurring distractors) are promising structural alternatives.

\textbf{Confidence abstention and downstream pipeline design.} Continuous-scoring judges produce a natural decision band that supports human-in-the-loop review. We did not implement and measure abstention strategies in the present work, but the score distributions suggest that a 10–15\% abstention rate could capture most of the residual judge errors after filtering. Quantifying the cost-benefit trade-off of abstention is a tractable empirical question.

Beyond these four directions, the broader research programme concerns the trustworthiness of automated evaluation in adversarial settings. As LLM-as-a-Judge becomes more widespread for production RAG evaluation, the need for evaluation infrastructure that is itself robust will only grow. The present thesis contributes a working pre-filter, a reusable dataset, an evaluation methodology, and an honest empirical characterisation of where current defences succeed and where they fail on a single multi-hop QA task. The larger problem - making evaluation pipelines as adversarially robust as the systems they evaluate, across the diverse settings in which RAG is deployed - remains open.

\section{Closing Remarks}

The Poisoned Evaluation Dataset, the pre-filter pipeline implementation, and the full experimental codebase are released alongside this thesis. The work was conducted under significant resource constraints, and we hope that the released artefacts make it easier for follow-up researchers to build on what has been demonstrated here without re-running the full experimental programme.

The empirical result of this work is uncomfortable. LLM judges are imperfect; corpus poisoning exploits their imperfections; a well-designed filter that removes the poison the dataset constructor injected does not, in the designs evaluated here, restore judge reliability - and may worsen it on the cases where reliability matters most. What this thesis adds is the rigorous experimental measurement of these claims across three contemporary judges, three attack types, and four noise levels on a single multi-hop QA task, alongside an analytical framework (the True/Survived decomposition) that distinguishes the cases where the filter succeeds from the cases where it does not. The single-task scope is a real limitation: the headline paradox should be read as established for HotpotQA-like multi-hop retrieval QA evaluation, with cross-task generalisation as the next experimental step rather than a finished claim.

The findings argue against complacency in deploying LLM-as-a-Judge in adversarial settings, and they argue for treating the judge as a first-class object of security analysis rather than as a downstream consumer that inherits its trustworthiness from upstream defences. The hope is that by establishing both the vulnerability and the limits of current mitigations on shared benchmark conditions, the work contributes to building automated evaluation pipelines that are demonstrably trustworthy in conditions where trust matters.
